I dette prosjektet analyseres lydklippet fra sangen “Seven Nation Army” 
av The White Stripes ved hjelp av Fast Fourier Transform (FFT) og Short-Time 
Fourier Transform (STFT). Målet er å undersøke hvordan Fourier-transformasjoner 
kan brukes til å identifisere og visualisere de dominerende 
frekvenskomponentene i et musikkstykke, samt forstå hvordan disse 
varierer over tid. FFT benyttes for å identifisere hovedfrekvensene og 
harmoniske overtoner i hele signalet, mens STFT muliggjør en tids-frekvensanalyse 
som viser dynamiske endringer i energi og rytmikk.

Analysen viser at de lavfrekvente komponentene mellom 50 - 120 Hz dominerer lydbildet, 
og representerer sangens karakteristiske basslinje. Gjennom spektrogrammet fremkommer 
tydelige mønstre som samsvarer med slagverk og rytmiske elementer i mellomfrekvensområdet 
(500 - 2000 Hz). Resultatene bekrefter at Fourier-transformasjoner gir et presist og 
intuitivt bilde av hvordan både harmoniske og rytmiske strukturer danner grunnlaget for musikken.

Videre drøftes effekten av ulike vindustyper og parametervalg i STFT-analysen. 
Hann-vinduet ble valgt for å balansere mellom tids- og frekvensoppløsning og 
for å redusere spektrallekkasje. Prosjektet viser hvordan matematiske transformasjoner 
kan anvendes i musikkteknologi for å analysere og tolke lydsignaler, og legger 
grunnlaget for videre arbeid med mer avanserte metoder som wavelet-transformasjon 
eller maskinlæringsbasert lydklassifisering.